\section{Outcome, scope, and relation to earlier work}
\label{sec:overview}

The project already has an exact recognizer. Its unsolved performance
problem is a worst-case complexity problem, not a missing correctness
criterion. In particular, the Khovanov backend decides triviality from the
rank of reduced homology over $\F$~\cite{kronheimer,shumakovitch}.
Kronheimer and Mrowka prove that a nontrivial knot has reduced rational
Khovanov rank greater than one. The universal coefficient theorem makes
the total reduced rank over $\F$ at least that rational rank, whereas the
unknot has reduced rank one. Over $\F$, unreduced homology is the direct
sum of two quantum-shifted copies of reduced homology, so dividing its
total rank by two gives the same decision criterion.
The present continuation concentrates
on making an expensive algebraic primitive substantially cheaper and on
providing a complementary, independently checkable obstruction.

\subsection{What this continuation contributes}

There are five main contributions. The first is an exact factorization of
the existing compiled cobordism plans. It applies to arbitrary morphisms in
the implemented algebra, including dense ones; it does not assume that a
particular knot family creates them. The second is an implementation whose
coefficient manipulation respects a bit-cost model. The distinction is
essential because a morphism with $m$ possible dots is a $2^m$-bit integer,
not an $m$-bit integer. The third recovers a grading invariant and uses it
to sharpen the sufficient parameter regime for exact component sharing.
The fourth is a finite-group certificate backend with
explicit construction and verification semantics. The fifth is a precise
separation of the local improvement from the global hypotheses that would
be needed for a quasi-polynomial recognizer.

\begin{center}
\small
\begin{tabular}{P{.26\linewidth}P{.40\linewidth}P{.23\linewidth}}
\toprule
Result & Mathematical content & Status \\
\midrule
Cobordism factorization & Input contraction, squarefree multiplication,
output reconstruction & Proved and implemented \\
Dense coefficient bound & $\poly(m)2^m$ bit cost for $2m$ endpoints &
Unconditional local bound \\
Homogeneous type count & $b^2\binom{m}{\floor{m/2}}$ in the sufficient
component-sharing exponent & Proved; conditional global use \\
Finite-group certificates & Noncyclic $A_5$ representation implies knottedness &
Proved; separate verifier \\
Wirtinger seed search & Exhaustive search with $q^{b-1}$ assignments for a
checked $b$-seed plan & Parameterized bound \\
Palette obstruction & $T(3,p)$ needs at least $p$ quandle colors &
Proved; no priority claim \\
Full quasi-polynomial recognition & Requires bounds on surviving data or a
constructive geometric hierarchy & Unresolved \\
\bottomrule
\end{tabular}
\end{center}

The subset-convolution algorithm itself is due to Bj\"orklund, Husfeldt,
Kaski, and Koivisto~\cite{bhkk}. Quandle coloring, normalization of one
meridian color, and the use of non-affine targets after a trivial Alexander
polynomial also precede this work~\cite{clark,fish}. Our claims concern the
factorization and implementation for this repository, the accompanying
bit-cost audit, the concrete certificate backend, and the deductions proved
below. We make no claim to have invented these underlying theories.

\Needspace{12\baselineskip}
\subsection{Sources and pinned baseline}

The computational baseline is the \texttt{fast/} subtree at ProveIt commit
\begin{center}
\texttt{ca61a1a5f967999d778d6efd7c882e77a2eb379d}.
\end{center}
Its tree identifier is
\texttt{b99a666269f660d932fc82ac8d1f29d3ae42231f}. The baseline source is
included in the archive so that comparisons do not depend on later
repository changes. Forty-one baseline unit tests passed in this
environment before modification.

The audit also read the structural-compression report now stored under
\texttt{reports/07}, and two recent continuation archives:
\nolinkurl{unknot_progress_20261007.zip} and
\nolinkurl{unknot_recognition_progress.zip}. These contain, respectively,
improved interlacement factorization, sparse Alexander elimination and a
pointed scanner; and an exact three-braid backend, Markov descent, and a
study of intrinsic scanner multiplicities. Report 07 contributes
homogeneous Seifert-graph certificates and exact component sharing. Those
results are antecedents, not results newly established by the present
measurements. The integration patch is against the pinned \texttt{fast/}
baseline, not an untested merge of every recent branch.

We inspected the README and manuscript catalogue of
\texttt{openai/math} at commit
\begin{center}
\texttt{adc7f1241b42e322a6451854ab7e4b4c146bf78a}.
\end{center}
The catalogue search did
not identify a directly applicable unknot-recognition theorem. No theorem
in the present argument is imported from that collection, and this limited
catalogue audit is not represented as a verification of its manuscripts.
The source record includes both repository revisions~\cite{proveit,openai}
and the relevant
primary literature.

\subsection{The precise quasi-polynomial target}

We use the standard broad terminology
\[
 T(n)=\exp\bigl((\log(n+2))^{O(1)}\bigr)
\]
for quasi-polynomial time. The repository's explicit objective
$n^{O(\log n)}$ is the stronger, particular scale
$\exp(O((\log n)^2))$.

Lackenby's detailed 2021 handout states an announced bound
$2^{O((\log n)^3)}$, equivalently $n^{O((\log n)^2)}$~\cite{lackenbytalk}.
His July 2026 preprint develops incompressible surfaces, boundary patterns,
and certifiable hierarchies~\cite{lackenby2026}. Its Section 9 separates
step counts from running time; several operations do not yet have the
precision needed for that cost analysis. The useful displayed search bound
is $L(g+1)^L$ in terms of hierarchy length $L$ and pattern complexity $g$.
The introduction reserves further speed-up for future work. Neither an
efficiently checkable hierarchy nor that parameterized visit count is, by
itself, a construction in the desired time bound.

Our end-to-end discussion therefore has two distinct questions. How much
does each local algebraic operation cost on the data actually present?
How large can the data become, and can the required data be constructed?
The first receives an unconditional improvement here. The second remains
the central research problem.

\subsection{Conventions and trust boundary}

Inputs are validated planar diagrams of a single knot. A diagram has $n$
crossings and integer edge labels; labels can be renumbered to have
$O(\log(n+2))$ bits. During a scan, $w=2m$ denotes the number of exposed
endpoints. Each object is a loopless matching on those endpoints, with a
homological degree. Closed circles are delooped. All cobordism
coefficients in this continuation are in $\F$; the integer or odd-prime
coefficient theories require additional sign and coefficient analysis.

We distinguish three kinds of evidence. A proof justifies an algorithm on
all inputs satisfying its hypotheses. An exhaustive finite test checks a
specified finite domain. A benchmark reports a measured execution on a
specified host and input. None substitutes for either of the others.
The accompanying proofs have not been formally checked in Lean, and the
computational checks are stated at their actual scope.
